\chapter{Computational Projection: KES, Spherepop, and Irreversible Event Calculi}
\label{ch:kes_spherepop}

\epigraph{The KES map formalizes the transition from current world-state to
hypothesis space as the basic cognitive operation. The question is always which
hypotheses remain admissible given what has already been committed.}{}

\section{The KES Map as Cognitive Architecture}

The Kinetic-Event Synthesis map $\Omega_t \to H_{t+1}$ formalizes the basic
cognitive operation: given a current configuration of constraints, evidence, and
prior commitments, what hypothesis spaces are admissible at the next step?

This is not prediction in the statistical sense. A system predicting the next
state must assign a probability distribution over all possible successors. A
system navigating a KES map must determine which successor hypothesis spaces
are admissible---which are consistent with the accumulated constraints---without
necessarily assigning probabilities to any of them. The difference is between
probabilistic modeling and admissibility filtering. Both are useful operations;
they are not the same operation.

\claimH The KES architecture is most useful as a heuristic framework for
characterizing cognitive systems that must select among competing hypotheses under
partial information and irreversibility constraints. It is not a complete
computational model but a formal language for describing the structure of the
selection problem.

\section{Spherepop Operators in Practice}

The four Spherepop operators constitute a minimal calculus for irreversible event
processing. Their practical interpretation:

\textbf{Pop} extracts committed events from possibility space. When a system
commits to an interpretation---a perceptual categorization, a theoretical claim,
a contractual obligation---it performs a Pop operation. The committed event
becomes actual; the alternatives do not disappear from the world but become
inaccessible to the committing system.

\textbf{Refuse} eliminates inadmissible trajectories. Before commitment, a system
may reject certain possibilities as inconsistent with its prior commitments,
its constraint set, or its evidence. Refuse is the active filtering operation
that produces the admissible trajectory space from the unconstrained space.

\textbf{Bind} establishes structural dependencies between events. When an event
$e_1$ is bound to an event $e_2$, the occurrence of $e_1$ constrains the
admissibility conditions for $e_2$. Bind creates the dependency graph that
constitutes the system's provenance structure.

\textbf{Collapse} commits a possibility space to a definite configuration. Where
Pop extracts a single event, Collapse resolves an entire horizon of possibilities
into a committed configuration. Collapse is irreversible: the alternatives that
were not selected are not merely deprioritized but structurally inaccessible
within the same system-instance.

\claimD All four operations are irreversible within a given system-instance. This
is a design commitment rather than a limitation: systems that can always reverse
their operations accumulate no structural residue and have no trajectory in the
meaningful sense.

\section{TARTAN and Trajectory-Aware Tiling}

\claimH The TARTAN framework---Trajectory-Aware Recursive Tiling with Annotated
Noise---provides a geometric interface between the abstract event calculus of
Spherepop and concrete computational substrate. Recursive tiling under trajectory
annotations preserves structural information that purely statistical approaches
discard.

The key feature is the annotated noise schema. Rather than discarding unresolved
structure as noise to be filtered, TARTAN archives it as annotated residue
available for future reinterpretation. This operationalizes the ``hallucination
as normal'' principle: the unresolved structure is not an error to be eliminated
but a resource to be managed. It constitutes exactly the non-trivial $H^1$ cocycles
of Chapter~\ref{ch:sheaf}: ambiguity stored for potential future repair.

\section{Gesture-First Symbolic Grounding}
\label{sec:gesture_grounding}

\claimP Sign language detection systems provide a particularly clean instantiation
of the admissibility framework because the symbolic and sub-symbolic layers are
architecturally separated and inspectable. A system trained to classify ASL gestures
using a convolutional detection pipeline operates in three distinct stages that
correspond directly to the vocabulary established in Chapter~\ref{ch:vocabulary}.

In the first stage, raw video frames arrive as unconstrained sensory input. Each
frame is a candidate world-state: a high-dimensional tensor encoding spatial and
temporal structure that the system has not yet committed to any interpretation.

In the second stage, a detection model applies learned spatial constraints to the
frame, generating confidence scores for each candidate sign class. This is
admissibility filtering: the model's learned weights encode the constraint set
$\mathcal{C}$ that determines which visual configurations are compatible with each
symbolic label. A confidence threshold (typically 0.5) operates as the formal
Refuse criterion---frames producing no detection above the threshold are excluded
from the symbolic trajectory without any label being committed.

In the third stage, detections above the threshold are committed via Pop: the
symbolic label enters the actuality set $\mathcal{E}$ and can propagate into
downstream systems (translators, interfaces, assistants).

\subsection{Gesture Before Symbol}

The architectural separation between detection and classification in these systems
formalizes a claim that runs across the companion work on gesture-first input theory:
that symbolic meaning is grounded in gestural trajectories rather than in static
symbolic configurations. The label ``hello'' does not exist independently in the
system's representational space; it exists only as a compression of a family of
admissible hand configurations, orientations, and spatial positions. The symbol is
a compressed attractor over a continuous gestural manifold.

\claimH This is a concrete computational implementation of Axiom~1 (Local Action
Only) from the Yarncrawler framework. The detection model operates locally: it
classifies the current frame without global access to conversational history,
intent, or semantic context. The symbolic output is generated from local spatial
constraints alone. Higher-order meaning---the pragmatic force of ``yes'' versus
``I love you'' in context---must be assembled by downstream layers that have access
to episodic and procedural context.

\subsection{The Label Set as Pragmatic Primitive}

The five signs chosen in a typical minimal ASL detection system---hello, yes, no,
thank you, I love you---are not chosen for semantic richness but for pragmatic
coverage. They are the primitives of social coordination and consent: acknowledgment,
affirmation, negation, gratitude, and affiliation. A system that can detect these
five signs can participate in the most basic structure of cooperative interaction,
even without any propositional content.

\claimP This connects to the admissibility geometry of social systems analyzed in
Chapter~\ref{ch:social_institutional}. The pragmatic primitives are exactly the
signs that regulate the admissibility of subsequent interaction: ``yes'' and ``no''
open and close trajectories; ``hello'' initializes the interaction state; ``thank
you'' signals trajectory completion; ``I love you'' establishes a binding that
constrains all subsequent exchange. The label set is a minimal Spherepop vocabulary
for social coordination.

\subsection{Transfer Learning as Boot Sequence}

The pipeline's use of transfer learning---initializing from SSD MobileNet weights
pretrained on COCO and fine-tuning on a small domain-specific dataset---is a
computational instance of the boot sequence concept from Chapter~\ref{ch:identity}.
The pretrained weights encode a lower-level constraint layer: the detection of
spatial features, edges, textures, and object-like configurations in natural images.
Fine-tuning on ASL gesture images adds a domain-specific constraint layer on top of
this foundation.

The fine-tuning layer is only meaningful because the pretrained layer has already
been initialized. A model trained from scratch on 75 gesture images (15 per class)
would fail to generalize; the same number of examples produces a functional
classifier when the lower layers have already been bootstrapped from millions of
natural images. The boot sequence is explicit in the configuration: the pipeline
specifies \texttt{fine\_tune\_checkpoint\_type = "detection"}, distinguishing the
pretrained lower layers (not updated) from the domain-specific upper layers (updated
from the small ASL dataset).

\begin{constraintboundary}
\textbf{Status:~\textbf{[P/H]}} The admissibility filtering and Spherepop
interpretations of the detection pipeline are phenomenological. The gesture-before-symbol
claim is a heuristic framework for understanding the grounding relationship,
not a derived result.

\textbf{Design implication:} Systems that ground symbolic output in gesture
trajectories (rather than in static image classification) naturally implement a
form of temporal qualification: the gesture unfolds over time, and the detection
system must integrate spatial constraints across frames to produce a stable
symbolic commitment. This is the temporal residue mechanism of Section~\ref{sec:marine}
applied to the sign recognition domain.

\textbf{Not claimed:} That the SSD MobileNet architecture is the correct
computational model of gestural grounding in biological cognition, or that the
five-sign label set is sufficient for full communicative participation.
\end{constraintboundary}

The reconstruction bound established in Chapter~\ref{ch:vocabulary}---that
$H(D \mid D') \geq N - M > 0$ for any compression with $M < N$---has a direct
computational consequence: any system operating under compression must interpolate
missing structure from prior regularities. In computational systems, this
interpolation produces hallucinations in the formal sense.

\claimD The design objective is therefore not zero hallucination but
\emph{bounded hallucination}: reconstruction errors that remain corrigible through
feedback rather than self-reinforcing through proxy stabilization. A system whose
hallucinations are corrigible is safe in a way that a system with zero apparent
hallucination but hidden proxy stabilization is not.

\claimC The uncertainty gradient finding from the Cheng, Broadbent, and Chappell
(2025) system is consistent with this framework: correct reasoning exhibits a
negative uncertainty gradient over time (hallucinations are detected and corrected),
while incorrect reasoning exhibits a positive or oscillating gradient (hallucinations
stabilize or escalate). The gradient is a corrigibility signal. A system monitoring
its own uncertainty gradient is implementing a weak version of the corrigibility
condition.

\begin{constraintboundary}
\textbf{Established:} KES as a formal language for admissibility-filtered hypothesis
selection. Spherepop operator semantics as a minimal irreversible event calculus.
The corrigibility condition as the correct design target for hallucination management.

\textbf{Heuristic~[H]:} TARTAN as a computational architecture. The annotated
noise schema is a heuristic engineering principle, not a formally derived result.

\textbf{Conjectural~[C]:} The uncertainty gradient as a universal corrigibility
signal. The specific claim requires empirical validation across diverse system types.

\textbf{Falsification:} A compressed inference system that achieves zero hallucination
under arbitrary compression ratios would falsify the reconstruction bound. A system
with positive uncertainty gradient that nonetheless converges on correct outputs
would weaken the corrigibility gradient claim.
\end{constraintboundary}

\section{Temporal Salience and Sparse Qualification}
\label{sec:marine}

The problem of salience detection under energetic and informational constraints
appears across biological cognition, signal processing, distributed sensing, and
adaptive computation. Conventional spectral methods frequently assume that
exhaustive decomposition is both feasible and desirable, yet biological systems
rarely possess the energetic or temporal capacity required for complete
reconstruction of an incoming signal field.

\claimH The following framework is introduced as a computational heuristic for
low-cost temporal qualification under bounded resources. The central claim is that
many systems can achieve robust environmental responsiveness through partial temporal
stabilization metrics rather than exhaustive representational completeness.

\subsection{Temporal Qualification Before Spectral Reconstruction}

Traditional signal analysis proceeds by transforming an incoming waveform into a
fully decomposed spectral representation. Biological systems, however, frequently
operate through earlier-stage temporal filtering mechanisms that identify local
instability, transient coherence, and salient directional discontinuities before
detailed reconstruction occurs.

Temporal qualification refers to the extraction of structurally relevant features
directly from evolving signal relations. Metrics such as local jitter, peak
persistence, envelope discontinuity, and transient regularity provide low-cost
approximations of environmental significance without requiring exhaustive
decomposition. This is admissibility pruning applied to signal processing: the
system progressively excludes trajectories inconsistent with local temporal
stability conditions rather than reconstructing every possible interpretation.

\subsection{Formal Specification}

Let an incoming signal be represented as a continuous trajectory
$x(t): \mathbb{R} \to \mathbb{R}^n$. Define a qualification window
$W_\tau(t) = [t-\tau, t+\tau]$ and a local temporal variation operator:
\[
  J_\tau(x,t) = \frac{1}{2\tau}\int_{t-\tau}^{t+\tau}
  \left\|\frac{dx}{ds}\right\| ds.
\]
The quantity $J_\tau$ measures local temporal instability. Signals with excessive
instability relative to environmental persistence thresholds are progressively
suppressed through local admissibility weighting:
\[
  A(x,t) = \exp(-\lambda J_\tau(x,t)),
\]
where $\lambda$ controls suppression strength.

Under a finite representational activation budget $\sum_{i=1}^N a_i(t) \leq
E_{\max}$, sparse qualification emerges through selective amplification of only
those trajectories satisfying admissibility conditions:
\[
  \mathcal{A}(t) = \bigl\{x_i \mid C(x_i,t) < \epsilon\bigr\},
\]
where $C(x_i,t)$ represents local incompatibility cost. Inference operates within
the admissible subset $\mathcal{A}(t) \subseteq \Omega(t)$ rather than across
the full trajectory space.

\subsection{Temporal Residue and Attentional Gating}

Temporal residue refers to the persistence of prior signal structure across
subsequent qualification intervals. Define recursive persistence:
\[
  R_n = \alpha R_{n-1} + (1-\alpha)\Phi(x_n),
\]
where $0 < \alpha < 1$ controls persistence memory and $\Phi(x_n)$ measures local
structural coherence. The attentional gate is then:
\[
  G(x_n) = \sigma(R_n - \theta),
\]
where $\sigma$ is a sigmoid function and $\theta$ is a gating threshold. Signals
maintaining coherent persistence accumulate larger residue and therefore increased
attentional weighting. Salience emerges from persistence under recursive
qualification rather than absolute magnitude alone.

\begin{constraintboundary}
\textbf{Status:~\textbf{[H]}} The framework proposes only that lightweight temporal
qualification heuristics may provide efficient approximations of salience under
bounded energetic conditions. Claims regarding subjective experience, emotional
authenticity, or universal cognitive equivalence are excluded from this formalism.

\textbf{Falsification:} A biologically plausible system that achieves equivalent
salience detection at lower computational cost without temporal qualification, or
a system where spectral exhaustion proves more efficient than temporal pruning
under realistic energetic constraints, would weaken the heuristic.
\end{constraintboundary}
